Let $e \in \mathcal{E}$ denote the managed entity under study. Depending on the experiment, $e$ may represent the aggregate cluster, an individual machine, or an application-level service identified by Alibaba's \texttt{app\_du}. For every entity, the time index is discretized to one-minute intervals. The observed demand signal is denoted by $x_{e,t} \in \mathbb{R}_{\ge 0}$.

\paragraph{Entity-level forecasting objective.}
Given a context window of length $L$, the forecasting layer receives
\[
\mathbf{x}_{e,t-L+1:t} = \{x_{e,t-L+1}, x_{e,t-L+2}, \dots, x_{e,t}\},
\]
and predicts future demand at horizons $\mathcal{H}=\{1,5,10\}$ minutes. For each $h \in \mathcal{H}$, the model produces at least two quantiles,
\[
\hat{x}^{(0.5)}_{e,t+h}, \qquad \hat{x}^{(0.9)}_{e,t+h},
\]
which correspond to the median and the upper decision quantile. The median supports conservative scale-in, while the upper quantile supports earlier scale-out under uncertainty.

\paragraph{Demand representation.}
The forecasting formulation is shared across all benchmarks, but the physical meaning of $x_{e,t}$ depends on the entity type. For aggregate and machine-level experiments, $x_{e,t}$ is the observed CPU-utilization signal after one-minute aggregation. For service-level experiments, $x_{e,t}$ is a CPU-core demand proxy obtained by aggregating container-level rows that belong to the same \texttt{app\_du} and summing
\[
\text{cpu\_request\_cores} \times \text{cpu\_utilization}.
\]
This keeps the forecasting route consistent while allowing the control layer to operate on service-relevant load units.

\paragraph{Required-capacity mapping.}
The controller acts on a required-capacity signal rather than on the raw forecast. For aggregate and machine-level experiments, the observed demand is first normalized to the policy capacity space and then converted into a minimum required capacity
\[
c^\star_{e,t} = \max\left((1+\rho)\tilde{x}_{e,t},\; c_{\min}\right),
\]
where $\rho$ is the headroom ratio, $\tilde{x}_{e,t}$ is normalized demand, and $c_{\min}$ is the minimum capacity floor.

For service-level experiments, the same logic is expressed in replica-equivalent units. The required replica-equivalent capacity becomes
\[
c^\star_{e,t} = \max\left(\frac{(1+\rho)x_{e,t}}{r_e},\; c_{\min}\right),
\]
where $r_e$ is the median per-container requested CPU cores for service $e$.
This mapping is important because it places reactive, lagged-tracking, and predictive service policies in the same decision space.

\paragraph{Reactive baseline.}
The reactive policy is threshold-based. If the current utilization exceeds an upper threshold, capacity is increased by at most one step. If utilization falls below a lower threshold, capacity is decreased by at most one step. This approximates the structure of common threshold-driven elasticity controllers and serves as the operational baseline.

\paragraph{Lagged target-tracking baseline.}
The second policy baseline is forecast-free but more explicit than threshold control. It drives the next action directly toward the previous-step required capacity:
\[
\begin{aligned}
c_{e,t+1} = \max\Bigl(&c_{\min},\\
&c_{e,t} + \mathrm{clip}(c^\star_{e,t} - c_{e,t}, -\Delta, \Delta)\Bigr).
\end{aligned}
\]
This policy uses the same minimum-capacity floor and maximum single-step change as the predictive controller, but it has no access to future uncertainty and therefore tests whether naive target chasing is already sufficient.

\paragraph{Predictive policy.}
The predictive controller consumes the two quantile forecasts after mapping them into the same required-capacity space as the reactive and lagged baselines. Scale-out decisions use the upper quantile together with an explicit scale-out margin, while scale-in decisions use the median together with a safety margin and a cooldown constraint:
{\small
\[
\begin{aligned}
c_{e,t+1} &= \min(c_{e,t} + \Delta, \hat{c}^{(0.9)}_{e,t+h}),
&& \text{if scale-out},\\
c_{e,t+1} &= \max(c_{e,t} - \Delta,
  (1+\gamma)\hat{c}^{(0.5)}_{e,t+h}),
&& \text{if scale-in},\\
c_{e,t+1} &= c_{e,t},
&& \text{otherwise}.
\end{aligned}
\]
}
where $\Delta$ is the maximum single-step change and $\gamma$ is the scale-in safety margin. Cooldown suppresses oscillatory scale-in without delaying urgent scale-out.

\paragraph{Forecasting backbone.}
The proposed model family is framed as an uncertainty-aware multi-scale temporal forecaster. In the present experiments, that interface is instantiated by \texttt{UA-MSTCN-Lite}, a paper-specific lightweight quantile-forest surrogate that emits $P50$ and $P90$ forecasts without introducing a heavyweight deep-learning dependency stack. \texttt{UA-MSTCN-Lite} is not the full deep UA-MSTCN architecture; it is the reproducible surrogate used to exercise the forecasting-to-control contract in this revision. The surrounding dataset contract, policy simulator, and evaluation protocol are kept stable so that a deeper temporal backbone can replace this surrogate in later work without changing the rest of the study design.

Concretely, \texttt{UA-MSTCN-Lite} uses a 60-minute history window and predicts the 1-, 5-, and 10-minute horizons. Its feature extractor includes lag values at 1, 2, 3, 5, 10, 15, 20, 30, 45, and 60 minutes; multi-scale windows of 3, 5, 10, 20, 30, and 60 minutes; per-window mean, standard deviation, minimum, maximum, median, 0.1 quantile, 0.9 quantile, endpoint difference, last-minus-mean, and slope; and additional short/long moving-average and volatility differences. Each horizon is fitted with a separate \texttt{RandomForestRegressor} using 200 trees, \texttt{min\_samples\_leaf=3}, \texttt{max\_features="sqrt"}, random seed $42+h$, and \texttt{n\_jobs=-1}. Per-tree predictions define the empirical $P50$ and $P90$ outputs. The $P90$ path is widened by the 0.99 quantile of positive training residuals relative to $P50$, while $P50$ and $P90$ are clipped to the training target support, with the upper path allowed up to 1.10 times the training maximum before the calibration floor. Multi-entity transfer experiments use the same implementation with pooled windows from the training entities.

\paragraph{Calibration objective.}
Because the controller consumes quantiles rather than only point forecasts, the evaluation also measures empirical coverage:
\[
\mathrm{cov}^{(q)}_{e,h} = \frac{1}{N}\sum_{i=1}^{N}\mathbb{I}\left(x_{e,t_i+h} \le \hat{x}^{(q)}_{e,t_i+h}\right),
\]
where $q \in \{0.5, 0.9\}$. Ideally, $P50$ and $P90$ coverage should remain close to 0.5 and 0.9, respectively. Interval width $\hat{x}^{(0.9)}_{e,t+h} - \hat{x}^{(0.5)}_{e,t+h}$ is tracked alongside coverage because overly wide intervals may protect SLA at the cost of resource efficiency.

\paragraph{Optimization targets.}
Forecasting quality is evaluated with MAE, RMSE, MAPE, and pinball-oriented quantile diagnostics. Control quality is evaluated with SLA violation rate, over-provisioning, under-provisioning, scaling-action count, and average reserved capacity. This two-layer evaluation is central to the paper because lower forecast error alone does not imply better autoscaling decisions.
